\documentclass[12pt,a4paper,twocolumn]{article}

\usepackage[a4paper,top=3cm,bottom=3cm,left=3cm,right=3cm]{geometry}
\setlength{\columnsep}{0.8cm} % Jarak antar kolom sesuai format 2 kolom TOR
\usepackage{mathptmx} % Times New Roman
\usepackage[indonesian]{babel}
\usepackage[none]{hyphenat}
\usepackage{setspace}
\usepackage[nopatch=footnote]{microtype}
\usepackage{anyfontsize}
\usepackage{booktabs}
\usepackage{tabularx}
\usepackage{amsmath}
\usepackage{titlesec}
\usepackage{caption}
\usepackage{footmisc}
\usepackage{hanging}
\usepackage{enumitem}
\setlist{itemsep=2pt, topsep=3pt, parsep=0pt}
\usepackage{xurl}
\usepackage[hidelinks]{hyperref}
\raggedbottom
\sloppy
\hbadness=10000

% Konfigurasi Penempatan Float 2-Kolom
\renewcommand{\topfraction}{0.9}
\renewcommand{\bottomfraction}{0.8}
\renewcommand{\textfraction}{0.05}
\renewcommand{\dbltopfraction}{0.95}
\renewcommand{\dblfloatpagefraction}{0.9}

% Format Judul Bagian
\titleformat{\section}{\bfseries\fontsize{11.5}{13.5}\selectfont}{\thesection.}{0.4em}{}
\titleformat{\subsection}{\bfseries\fontsize{11}{12.5}\selectfont}{\thesubsection.}{0.4em}{}
\titleformat{\subsubsection}{\itshape\fontsize{10.5}{12}\selectfont}{\thesubsubsection.}{0.4em}{}
\titlespacing*{\section}{0pt}{4pt plus 1pt minus 1pt}{1.5pt plus 1pt minus 1pt}
\titlespacing*{\subsection}{0pt}{3.5pt plus 1pt minus 1pt}{1.2pt plus 1pt minus 1pt}
\titlespacing*{\subsubsection}{0pt}{3pt plus 1pt minus 1pt}{1pt plus 1pt minus 1pt}

% Pengaturan Paragraf dan Spasi
\setlength{\parindent}{0.65cm}
\setlength{\parskip}{0pt}
\setstretch{0.98} % Proporsi line spacing rapat presisi
\setlength{\footnotesep}{4pt}
\setlength{\skip\footins}{6pt plus 2pt minus 2pt}

% Format Caption Tabel
\captionsetup[table]{labelsep=period, justification=centering, font={bf,small}, skip=4pt}

\begin{document}

\twocolumn[
    \begin{@twocolumnfalse}
        % =========================================================================
        % JUDUL NASKAH
        % =========================================================================
        \begin{center}
            {\fontsize{13.5}{16}\selectfont \textbf{\textsc{Menata Ulang Relasi Yudikatif-Eksekutif-Legislatif: Urgensi Otonomi Anggaran Lembaga Yudikatif dan Pemisahan Rezim Manajemen Aparatur Peradilan}}\par}
            \vspace{0.4em}
            {\fontsize{10.5}{12}\selectfont \textbf{MOCH. ARDANY CHABIB}\par}
            {\fontsize{9.5}{11}\selectfont \textbf{PRAKTISI HUKUM (PANITERA PENGGANTI)}\par}
            {\fontsize{9}{10.5}\selectfont Email: ardanych@gmail.com\par}
            \vspace{0.4em}

            \begin{minipage}{0.93\textwidth}
                \begin{center}
                    \textbf{\fontsize{10}{12}\selectfont Abstrak}
                \end{center}
                \vspace{-0.3em}
                \fontsize{9.5}{11.8}\selectfont
                Kebijakan Satu Atap (\textit{one-roof system}) pascareformasi belum mewujudkan independensi peradilan sejati akibat subordinasi anggaran di bawah eksekutif dan penundukan kepegawaian dalam rezim aparatur sipil negara. Kajian yuridis-normatif berpendekatan kelembagaan baru ini bertujuan merekonstruksi relasi horizontal antarlembaga negara melalui pelembagaan otonomi anggaran yudikatif dan pemisahan korps Aparatur Sipil Peradilan (\textit{judicial civil service}). Hasil kajian mengungkap bahwa ketiadaan otonomi fiskal mendegradasi martabat peradilan menjadi pemohon anggaran di bawah kementerian eksekutif. Relasi penganggaran harus ditransformasikan menjadi negosiasi konstitusional langsung antara Mahkamah Agung dan DPR RI selaku pemegang hak bujet (\textit{power of the purse}) melalui kombinasi alokasi bertahap menuju 1,0\%--1,5\% APBN (\textit{statutory floor}), pengajuan mandiri (\textit{direct submission}), dan larangan pemotongan anggaran (\textit{non-diminution rule}). Alokasi mandiri ini mendesak guna membiayai belanja modal lebih dari 400 gedung pengadilan yang belum memenuhi standar prototipe (terutama pada satuan kerja daerah pemekaran dan wilayah 3T), memulihkan kesejahteraan aparatur peradilan, serta menuntaskan krisis tunggakan ribuan eksekusi perkara perdata nasional melalui penyediaan pos biaya bantuan pengamanan kepolisian, pengukuran BPN, penilai harta (KJPP), dan lelang. Selain itu, pemaknaan reduksionis eksekutif yang memusatkan kekuasaan kehakiman semata pada figur hakim—sebagaimana kenaikan tunjangan jabatan via PP No. 42 Tahun 2025 tanpa kenaikan bagi panitera dan jurusita—memicu disparitas kesejahteraan internal yang ekstrem. Penyeragaman birokrasi ASN melahirkan kebuntuan nomenklatur jabatan fungsional panitera dan jurusita di KemenPAN-RB serta ketergantungan teknis administrasi kepegawaian (termasuk kepangkatan dan golongan ruang hakim selaku pejabat negara) pada BKN, sekaligus mematikan tangga karier yudisial, di mana batas kaku usia pengangkatan hakim maksimal 40 tahun dalam tiga UU paket peradilan 2009 tetap menghanguskan peluang Panitera Pengganti senior untuk beralih menjadi hakim karier, meskipun PERMA No. 5 Tahun 2025 telah membuka jalur pengadaan dari PNS Mahkamah Agung. Menegakkan prinsip bahwa keadilan bukan ajang coba-coba teori sekaligus mengatasi defisit majelis hakim pada satuan kerja berbeban perkara tinggi, pola pembinaan hakim wajib ditata bertingkat melalui residensi yudisial sebagai anggota majelis di Pengadilan Kelas IA dan Kelas IA Khusus sebelum memimpin persidangan mandiri. Analog dengan pemisahan rezim TNI dan Polri, kemandirian yudikatif menuntut pelembagaan korps aparatur di bawah kewenangan eksklusif Mahkamah Agung yang diawasi sistem merit independen. Badan Pengkajian MPR RI direkomendasikan mengartikulasikan haluan ini ke dalam Pokok-Pokok Haluan Negara (PPHN) yang ditindaklanjuti melalui revisi UU Kekuasaan Kehakiman berorientasi sistem organisasi pengadilan, harmonisasi tiga UU Peradilan, pengesahan RUU Jabatan Hakim, dan revisi UU Keuangan Negara.

                \vspace{0.4em}
                \noindent \textbf{Kata Kunci:} Aparatur Sipil Peradilan; Kemerdekaan Peradilan; Otonomi Anggaran; Pokok-Pokok Haluan Negara; Relasi Antarlembaga Negara.
            \end{minipage}
            \vspace{0.5em}
        \end{center}
    \end{@twocolumnfalse}
]

% =========================================================================
% 1. PENDAHULUAN
% =========================================================================
\section{Pendahuluan}

Pasal 24 ayat (1) UUD NRI Tahun 1945 menjamin kemerdekaan kekuasaan kehakiman sebagai pilar negara hukum demokratis untuk mengakhiri subordinasi peradilan era Orde Baru di bawah Departemen Kehakiman.\footnote{Jimly Asshiddiqie, \textit{Konstitusi dan Konstitusionalisme Indonesia}, Edisi Kedua (Jakarta: Sinar Grafika, 2021), hlm. 152; Sebastiaan Pompe, \textit{The Indonesian Supreme Court: A Study of Institutional Collapse and Reconstruction} (Ithaca: Southeast Asia Program Publications, Cornell University, 2005), hlm. 88--92.} Pascareformasi, sistem satu atap (\textit{one-roof system}) melalui UU No. 35 Tahun 1999, UU No. 4 Tahun 2004, dan UU No. 48 Tahun 2009 mengalihkan aspek organisasi, administrasi, dan finansial empat lingkungan peradilan ke Mahkamah Agung.\footnote{Bagir Manan, \textit{Menegakkan Hukum Suatu Pencarian} (Jakarta: Asosiasi Advokat Indonesia, 2009), hlm. 45--47.} Namun, lebih dari dua dekade berjalan, penyatuan ini terjebak paradoks: peradilan merdeka dalam memutus perkara (\textit{decisional independence}), tetapi kemandirian institusionalnya masih tersandera instrumen anggaran dan kepegawaian eksekutif.\footnote{Zainal Arifin Mochtar, \textit{Lembaga Negara Independen: Dinamika Perkembangan dan Urgensinya dalam Struktur Ketatanegaraan Indonesia} (Depok: Rajawali Pers, 2021), hlm. 112.}

Kerentanan tersebut tampak pada dua aspek krusial. Pertama, ketiadaan otonomi anggaran: berdasarkan UU No. 17 Tahun 2003 tentang Keuangan Negara, Mahkamah Agung disamakan dengan kementerian teknis sehingga pagu dan belanja operasionalnya disaring sepihak oleh Kementerian Keuangan dan Bappenas. Kedua, distorsi rezim aparatur: kendati Pasal 19 UU No. 48/2009, Pasal 57--58 UU No. 20/2023 tentang ASN, dan Putusan MK No. 32/PUU-XII/2014 mendudukkan hakim sebagai pejabat negara,\footnote{Lihat Pasal 19 Undang-Undang Nomor 48 Tahun 2009 tentang Kekuasaan Kehakiman; Pasal 57 dan Pasal 58 Undang-Undang Nomor 20 Tahun 2023 tentang Aparatur Sipil Negara; serta Putusan Mahkamah Konstitusi Nomor 32/PUU-XII/2014 perihal Pengujian UU Nomor 5 Tahun 2014 tentang ASN terhadap UUD 1945, diucapkan tanggal 20 April 2015, bagian Pertimbangan Hukum.} kenaikan tunjangan hakim via PP No. 42 Tahun 2025 tanpa perbaikan sepadan bagi panitera dan jurusita justru memicu disparitas ekstrem, sementara penyeragaman pegawai non-hakim di bawah KemenPAN-RB dan BKN menimbulkan kebuntuan nomenklatur dan loyalitas ganda.

Asimetri relasi ini mengubah prinsip \textit{checks and balances} menjadi koersi institusional tidak langsung terhadap cabang yudikatif, sehingga penataannya menjadi agenda mendesak bagi MPR RI dalam merumuskan Pokok-Pokok Haluan Negara (PPHN). Kajian ini menjawab tiga permasalahan: (1) implikasi subordinasi anggaran dan aparatur terhadap kemerdekaan peradilan; (2) perbandingan internasional serta rancang bangun otonomi anggaran dan pemisahan rezim aparatur peradilan; dan (3) rekomendasi operasional bagi PPHN serta legislasi nasional.

Penelitian ini bertujuan memformulasikan desain otonomi anggaran yudikatif dan pemisahan korps Aparatur Sipil Peradilan (\textit{Judicial Civil Service}) bagi Badan Pengkajian MPR RI. Mengingat batasan ruang lingkup format makalah ilmiah, kajian ini menghasilkan sebuah \textbf{hipotesis institusional dan model kebijakan konseptual} yang memerlukan pendalaman lanjutan serta simulasi aktuaria fiskal bersama MPR RI dan kementerian terkait.

% =========================================================================
% 2. METODE ATAU PENDEKATAN KAJIAN
% =========================================================================
\section{Metode atau Pendekatan Kajian}

Kajian ini menggunakan penelitian hukum yuridis-normatif (\textit{doctrinal legal research}) berpendekatan kelembagaan baru (\textit{new institutionalism}) dalam hukum tata negara.\footnote{Peter Mahmud Marzuki, \textit{Penelitian Hukum}, Edisi Revisi (Jakarta: Kencana Prenada Media Group, 2017), hlm. 35.} Penelitian menelaah keselarasan vertikal dan horizontal norma Pasal 24 UUD NRI Tahun 1945 dengan perundang-undangan sektor keuangan dan kepegawaian.

Pendekatan meliputi: (1) \textit{Statute Approach}, meneliti paket perundang-undangan peradilan, keuangan negara, ASN, serta putusan MK terkait; (2) \textit{Conceptual Approach}, membedah doktrin pemisahan kekuasaan, asas kemerdekaan peradilan, teori otonomi anggaran, dan \textit{judicial civil service}; serta (3) \textit{Comparative Approach}, menelaah model Kosta Rika (alokasi minimum konstitusional), Filipina (otonomi fiskal \textit{non-diminution}), dan AS (\textit{direct budget submission}).

Bahan hukum primer dan sekunder dianalisis secara kualitatif preskriptif guna menghasilkan argumentasi hukum kokoh dan rekomendasi kebijakan aplikatif bagi Badan Pengkajian MPR RI.

% =========================================================================
% 3. ANALISIS DAN PEMBAHASAN
% =========================================================================
\section{Analisis dan Pembahasan}

\subsection{Landasan Konstitusional dan Teoretis Kemerdekaan Kekuasaan Kehakiman: Meninjau Ulang Checks and Balances Horizontal}

Secara klasik, doktrin pemisahan kekuasaan dari Montesquieu menuntut kemandirian absolut cabang yudikatif dari legislatif dan eksekutif,\footnote{Montesquieu, \textit{The Spirit of Laws}, diterjemahkan oleh Thomas Nugent (New York: Colonial Press, 1900), hlm. 151--153; lihat pula Jacob T. Levy, ``The Separation of Powers,'' dalam Richard Bellamy dan Jeff King (eds.), \textit{The Cambridge Handbook of Constitutional Theory} (Cambridge: Cambridge University Press, 2025), hlm. 279--296.} namun Alexander Hamilton dalam \textit{Federalist No. 78} mengingatkan bahwa peradilan tidak memegang kekuatan pemaksa (\textit{the sword}) maupun kendali anggaran (\textit{the purse}), sehingga menjadi cabang paling rentan di hadapan agresi dua cabang politik lainnya.\footnote{Alexander Hamilton, James Madison, dan John Jay, \textit{The Federalist Papers}, diedit oleh Clinton Rossiter (New York: New American Library, 1961), No. 78, hlm. 465.}

Dalam teori hukum ketatanegaraan modern, konsep kemerdekaan kekuasaan kehakiman tidak boleh hanya dipahami secara sempit sebagai kebebasan hakim secara individual dalam menjatuhkan putusan di ruang sidang tanpa intervensi (\textit{substantive/decisional independence}). Shimon Shetreet dan Sophie Turenne menegaskan bahwa independensi peradilan bertumpu pada empat dimensi terpadu: (1) \textit{substantive independence}; (2) \textit{personal independence}; (3) \textit{collective/internal independence}; dan (4) \textit{institutional or administrative independence}.\footnote{Shimon Shetreet dan Sophie Turenne, \textit{Judges on Trial: A Study of the Appointment and Accountability of the English Judiciary}, Edisi Kedua (Cambridge: Cambridge University Press, 2013), hlm. 68--75.} Dimensi keempat mensyaratkan bahwa lembaga peradilan wajib memiliki kemandirian mutlak dari pengaruh cabang eksekutif dan legislatif dalam mengelola administrasi peradilan, mengelola personil pendukung, serta menentukan dan membelanjakan anggarannya sendiri.

Dalam konteks hukum tata negara Indonesia, Pasal 24 ayat (1) UUD NRI Tahun 1945 mengkonstruksikan kemerdekaan kekuasaan kehakiman dalam kedudukan horizontal yang setara dan saling mengimbangi (\textit{checks and balances}) dengan Presiden selaku pemegang kekuasaan pemerintahan (Pasal 4 ayat (1)) dan Dewan Perwakilan Rakyat selaku pemegang kekuasaan membentuk undang-undang dan hak anggaran (Pasal 20 ayat (1) dan Pasal 20A ayat (1)).\footnote{Saldi Isra, \textit{Pergeseran Kekuasaan Legislatif: Mengurai Hubungan Presiden dan DPR Pasca Perubahan UUD 1945} (Jakarta: Rajawali Pers, 2020), hlm. 165--172; lihat pula Restu Permadi dan Fifiana Wisnaeni, ``Tinjauan Hukum Kemandirian dan Independensi Mahkamah Agung di dalam Sistem Ketatanegaraan Indonesia,'' \textit{Jurnal Pembangunan Hukum Indonesia} 2, no. 3 (2020): 399--415, \url{https://doi.org/10.14710/jphi.v2i3.399-415}.} Namun, prinsip kesetaraan horizontal ini mengalami degradasi mendasar dalam operasionalisasinya. Kekuasaan kehakiman berada dalam posisi inferior karena cabang eksekutif menguasai seluruh instrumen regulasi birokrasi dan fiskal negara, sementara cabang legislatif menguasai persetujuan penganggaran.

Ketika relasi antarlembaga negara berjalan asimetris, mekanisme \textit{checks and balances} berubah menjadi celah koersi institusional. Jika cabang eksekutif atau legislatif merasa terancam oleh putusan pengadilan yang membatalkan undang-undang atau kebijakan pemerintah, mereka dapat melakukan pembalasan tidak langsung (\textit{judicial retaliation}) melalui instrumen pemotongan anggaran, pembatasan formasi kepegawaian, atau penundaan hak-hak keuangan aparatur peradilan.\footnote{Rosalind Dixon dan David Landau, \textit{Abusive Constitutional Borrowing: Legal Globalization and the Subversion of Liberal Democracy} (Oxford: Oxford University Press, 2021), hlm. 114--118; lihat pula James W. Douglas dan Roger E. Hartley, ``The Politics of Court Budgeting in the States: Is Judicial Independence Threatened by the Budgetary Process?'' \textit{Public Administration Review} 63, no. 4 (2003): 441--454, \url{https://doi.org/10.1111/1540-6210.00307}.} Dengan demikian, independensi yudikatif yang dijamin konstitusi menjadi ilusi jika tidak dibentengi oleh kemandirian fiskal dan otonomi aparatur.

Gagasan otonomi anggaran dan aparatur ini tidak boleh disalahpahami sebagai upaya membangun kekuasaan kehakiman absolut tanpa kontrol (\textit{imperial judiciary}). Kemandirian kelembagaan wajib berjalan seiring dengan akuntabilitas publik agar kekuasaan kehakiman tidak terjebak menjadi kekuasaan yang tak terkendali, sebab dalam kerangka tata kelola peradilan modern, independensi dan akuntabilitas merupakan dua dimensi yang saling menguatkan dan dapat diukur secara objektif—termasuk pada transparansi pengelolaan anggaran dan administrasi pengadilan.\footnote{Rahayu Prasetianingsih, ``Akuntabilitas Kekuasaan Kehakiman,'' \textit{Jurnal Konstitusi} 8, no. 5 (2011): 829--848, \url{https://doi.org/10.31078/jk858}; Frans van Dijk dan Geoffrey Vos, ``A Method for Assessment of the Independence and Accountability of the Judiciary,'' \textit{International Journal for Court Administration} 9, no. 3 (2018): 1--21, \url{https://doi.org/10.18352/ijca.276}.} Dalam paradigma tata kelola fiskal modern, kekuasaan kehakiman tetap terikat pada standar transparansi keuangan negara: tunduk pada audit kepatuhan dan kinerja Badan Pemeriksa Keuangan (BPK), mempublikasikan Laporan Keuangan Pemerintah Pusat secara transparan, serta mempertanggungjawabkan pemanfaatan anggaran dalam rapat dengar pendapat terbuka di hadapan DPR RI. Demikian pula dalam dimensi kepegawaian, pengelolaan eksklusif korps \textit{Judicial Civil Service} oleh Mahkamah Agung tetap wajib diawasi standar meritokrasinya secara eksternal oleh lembaga pengawas sistem merit yang independen—selaras dengan semangat Putusan Mahkamah Konstitusi Nomor 121/PUU-XXII/2024\footnote{Mahkamah Konstitusi RI, \textit{Putusan Nomor 121/PUU-XXII/2024}, perihal Pengujian Undang-Undang Nomor 20 Tahun 2023 tentang Aparatur Sipil Negara terhadap UUD 1945, yang memerintahkan pembentukan lembaga pengawas sistem merit independen dalam tenggat waktu dua tahun.}—guna mencegah risiko nepotisme internal. Otonomi fiskal dan aparatur yudikatif semata-mata menyingkirkan campur tangan diskresioner kementerian eksekutif dalam menyaring pagu, memotong anggaran, atau mendikte formasi peradilan secara sepihak, bukan melepaskan peradilan dari pengawasan konstitusional.

\subsection{Paradoks Sistem Satu Atap: Subordinasi Fiskal dan Belenggu Power of the Purse}

Kelahiran kebijakan Satu Atap (\textit{one-roof system}) pada tahun 1999 merupakan lompatan sejarah dalam melepaskan badan peradilan dari cengkeraman Departemen Kehakiman. Namun, pemindahan payung institusional tersebut terbukti belum menyentuh jantung persoalan tata kelola keuangan negara. Berdasarkan Undang-Undang Nomor 17 Tahun 2003 tentang Keuangan Negara dan Undang-Undang Nomor 1 Tahun 2004 tentang Perbendaharaan Negara, kedudukan Mahkamah Agung disamakan dengan kementerian negara biasa (Pengguna Anggaran pada Bagian Anggaran 005).\footnote{Undang-Undang Nomor 17 Tahun 2003 tentang Keuangan Negara, Lembaran Negara RI Tahun 2003 Nomor 47, Pasal 6--8.} Kajian doktrinal mutakhir mengonfirmasi bahwa kewenangan Menteri Keuangan mengesahkan DIPA lembaga peradilan berdasarkan Pasal 7 ayat (2) huruf b UU Nomor 1 Tahun 2004 telah mereduksi ketentuan ``mata anggaran tersendiri'' menjadi sekadar label administratif (\textit{executive encroachment}), sejalan dengan temuan empiris politik penganggaran pengadilan yang menunjukkan kerentanan independensi yudikatif akibat ancaman pemotongan dan penyaringan anggaran oleh cabang politik.\footnote{Fahri Bachmid dan Diani Indah Rachmitasari, ``Rekonstruksi Kemandirian Anggaran Kekuasaan Kehakiman dalam Sistem Ketatanegaraan Indonesia,'' \textit{Jurnal Legislasi Indonesia} 23, no. 2 (2026): 263--278, \url{https://doi.org/10.54629/jli.v23i2.1733}; James W. Douglas dan Roger E. Hartley, op. cit., hlm. 441--454.}

Penyamaan status ini menimbulkan distorsi yang sangat parah dalam tiga tahapan siklus anggaran:
\begin{enumerate}
    \item \textbf{Perencanaan Pagu:} Usulan anggaran Mahkamah Agung tidak diajukan langsung ke DPR, melainkan disaring sepihak oleh Kementerian Keuangan dan Bappenas melalui Pagu Indikatif dan Alokasi Anggaran.
    \item \textbf{Pembahasan Parlemen:} Ketiadaan saluran pengajuan langsung memaksa pimpinan peradilan melakukan lobi informal di Komisi III DPR RI demi memperjuangkan anggaran operasional. Praktik lobi anggaran personal ini rentan menciptakan kompromi politik yang mencederai etika dan marwah peradilan.
    \item \textbf{Pelaksanaan Anggaran:} Mahkamah Agung terikat aturan perbendaharaan eksekutif (mekanisme DIPA dan revisi anggaran Kemenkeu), bahkan terkena pemotongan anggaran (\textit{automatic adjustment}) dan kebijakan efisiensi via Inpres tanpa konsultasi.\footnote{Christina Clarissa Intania, \textit{Policy Assessment: Dampak Kebijakan Efisiensi Anggaran terhadap Independensi Peradilan} (Jakarta: The Indonesian Institute, 2025), hlm. 4--10.}
\end{enumerate}

Untuk membuktikan secara empiris kerentanan fiskal tersebut, Tabel \ref{tab:anggaran_ma} memperlihatkan perkembangan usulan atau pagu awal Mahkamah Agung beserta rasionya terhadap APBN dibandingkan dengan pagu yang disetujui (pagu akhir DIPA) periode 2020--2027.

\begin{table*}[t]
    \centering
    \caption{Data Usulan, Pagu Disetujui, dan Rasio Anggaran Mahkamah Agung RI terhadap APBN (TA 2020--2027)}
    \label{tab:anggaran_ma}
    \footnotesize
    \begin{tabularx}{\textwidth}{@{} c >{\centering\arraybackslash}p{2.2cm} >{\centering\arraybackslash}p{3.1cm} >{\centering\arraybackslash}p{3.1cm} >{\raggedright\arraybackslash}X @{}}
        \toprule
        \textbf{TA} & \textbf{Belanja APBN} & \textbf{Usulan / Pagu Awal MA} & \textbf{Disetujui / Pagu Akhir MA} & \textbf{Keterangan} \\
        & \textbf{(Triliun Rp)} & \textbf{(Triliun Rp \& \% APBN)} & \textbf{(Triliun Rp \& \% APBN)} & \\
        \midrule
        2020 & 2.589,90 & 10,59 (0,409\%) & 9,85 (0,380\%) & Realisasi Rp9,40 T \\
        2021 & 2.750,00 & 11,24 (0,409\%) & 10,73 (0,390\%) & Realisasi Rp10,51 T \\
        2022 & 2.714,20\textsuperscript{a} & 12,518 (0,461\%) & 11,514 (0,371\%) & \textit{Refocusing} Rp983,84 M; realisasi Rp11,28 T \\
        2023 & 3.061,20 & 12,260 (0,400\%) & 11,911 (0,389\%) & Pagu DIPA TA 2023; realisasi Rp11,59 T \\
        2024 & 3.325,10 & 11,530 (0,347\%) & 11,925 (0,359\%) & Pagu indikatif awal; realisasi Rp11,60 T \\
        2025 & 3.621,30 & 15,160 (0,419\%)\textsuperscript{b} & 13,149 (0,363\%) & Pagu awal Rp12,68 T; realisasi Rp12,90 T \\
        2026 & 3.842,70 & --- & 22,957 (0,597\%) & Lonjakan +74,6\% YoY (tunjangan hakim) \\
        2027 & 4.106,26 & 27,253 (0,664\%)\textsuperscript{c} & 23,340 (0,568\%) & RAPBN Rp19,341 T (0,471\%); final Okt '26 \\
        \bottomrule
    \end{tabularx}
    \vspace{4pt}
    \begin{minipage}{\textwidth}
        \footnotesize \textit{Catatan: \textsuperscript{a}Pada TA 2022, rasio awal (0,461\%) dihitung terhadap APBN awal Rp2.714,20 T, sedangkan rasio revisi (0,371\%) dihitung terhadap APBN revisi Perpres 98/2022 Rp3.106,40 T (\textit{refocusing} MA Rp983,84 M). \textsuperscript{b}Pagu indikatif Rp12,15 T + usulan tambahan Rp3,01 T. \textsuperscript{c}Pagu indikatif Rp16,95 T + usulan tambahan Rp10,303 T (RAPBN Rp19,341 T atau 71,0\% dari usulan). Sumber: Diolah dari Laporan Tahunan MA RI (2020--2024), Laporan Keuangan MA Audited TA 2025, Rekapitulasi Resmi MA RI TA 2026--2027, dan Informasi APBN Kemenkeu RI (2022--2025).}\footnote{Mahkamah Agung RI, \textit{Laporan Tahunan Mahkamah Agung RI Tahun 2020} (Jakarta: MA RI, 2021); \textit{Laporan Tahunan Tahun 2021} (Jakarta: MA RI, 2022); \textit{Laporan Tahunan Tahun 2022} (Jakarta: MA RI, 2023); \textit{Laporan Tahunan Tahun 2023} (Jakarta: MA RI, 2024), hlm. 128--135; \textit{Laporan Tahunan Tahun 2024} (Jakarta: MA RI, 2025), hlm. 94--98; \textit{Laporan Keuangan MA RI Audited TA 2025} (Jakarta: MA RI, 2026); Biro Hukum dan Humas MA RI, ``Plt. Sekretaris MA Dorong Penguatan Infrastruktur Peradilan pada Tahun Anggaran 2027,'' \textit{mahkamahagung.go.id}, 2026; serta Kementerian Keuangan RI, \textit{Informasi APBN Tahun Anggaran 2022--2025} (Jakarta: Ditjen Anggaran Kemenkeu RI, 2022--2025).}
    \end{minipage}
\end{table*}

Data empiris pada Tabel \ref{tab:anggaran_ma} mengungkap dua fase krusial tata kelola fiskal kekuasaan kehakiman. Pertama, \textbf{periode stagnasi (2022--2025)}, ketika nominal pagu MA bergerak landai di rentang Rp11,51 triliun hingga Rp13,15 triliun sementara belanja APBN terus ekspansif, sehingga rasio anggaran MA menyentuh titik terendah pada TA 2024 sebesar 0,359\% dan TA 2025 sebesar 0,363\%. Pada fase ini, usulan kebutuhan riil MA kerap mengalami pemangkasan signifikan: dalam Rapat Kerja Komisi III DPR RI tanggal 13 Juni 2024, MA memaparkan pagu indikatif TA 2025 sebesar Rp12,15 triliun dan memperjuangkan usulan tambahan Rp3,01 triliun (total usulan Rp15,16 triliun), namun pagu awal APBN 2025 hanya menyetujui Rp12,68 triliun (83,6\% dari usulan) sebelum naik menjadi Rp13,149 triliun (86,7\%) pada pagu akhir. Dari alokasi terbatas tersebut, belanja pegawai rutin menyerap rata-rata 65,1\% hingga 68,8\% pagu akhir (mencapai Rp9,05 triliun pada TA 2025 \textit{audited}), menyisakan ruang belanja modal sangat sempit sekitar 10,0\% (Rp1,32 triliun). Kedua, \textbf{fase lonjakan struktural (2026--2027)}, ketika pagu MA melonjak 74,6\% menjadi Rp22,957 triliun (0,597\% APBN) pada TA 2026 akibat penyesuaian tunjangan hakim pasca-PP No. 42/2025. Namun, celah antara usulan MA dan pagu yang disetujui tetap berulang pada TA 2027: dari pagu indikatif awal Rp16,95 triliun ditambah usulan tambahan Rp10,303 triliun (total kebutuhan riil Rp27,253 triliun, mencakup Rp5,28 triliun belanja modal infrastruktur), RAPBN awalnya hanya mengakomodasi Rp19,341 triliun (71,0\% dari usulan; 0,471\% APBN) sebelum disesuaikan dalam RKA Oktober 2026 menjadi Rp23,340 triliun (85,6\% dari usulan; 0,568\% APBN). Kondisi ini menegaskan bahwa dari total 930 satuan kerja pengadilan secara nasional, pemenuhan lebih dari 400 gedung pengadilan belum standar prototipe tetap tersendat oleh mekanisme negosiasi pagu tahunan di bawah kendali eksekutif.

Ketiadaan otonomi fiskal ini melahirkan anomali etika dan kelembagaan yang sistemik di tingkat daerah. Akibat minimnya pagu belanja modal dalam DIPA, satuan kerja pengadilan di pelosok daerah terpaksa mengandalkan bantuan hibah tanah, pinjam pakai gedung, hingga kendaraan operasional dari Pemerintah Daerah (Bupati/Walikota/Gubernur). Fenomena ``mengemis fasilitas'' kepada eksekutif daerah ini melegalkan konflik kepentingan struktural (\textit{structural conflict of interest}), mengingat kepala daerah dan aparatur dinasnya merupakan subjek yang berpotensi besar menjadi pihak berperkara dalam kasus tindak pidana korupsi, sengketa perdata aset, maupun sengketa peradilan tata usaha negara.

Dampak paling destruktif dari keterbatasan anggaran peradilan adalah \textbf{menumpuknya ribuan tunggakan eksekusi perkara perdata di seluruh Indonesia} (\textit{civil judgment execution backlog}). Dalam praktik hukum acara perdata, beban biaya eksekusi sepenuhnya dibebankan kepada pemohon melalui panjar biaya perkara yang sangat mahal. Ketika pemohon tidak memiliki kemampuan finansial atau ketika eksekusi menghadapi perlawanan fisik di lapangan, pelaksanaan eksekusi putusan yang telah berkekuatan hukum tetap (\textit{inkracht van gewijsde}) terhenti bertahun-tahun. Hal ini berpangkal dari ketiadaan mata anggaran negara untuk mendanai instrumen dan institusi pendukung eksternal eksekusi perdata, yang mencakup:
\begin{enumerate}
    \item \textbf{Biaya Bantuan Pengamanan Kepolisian (PAM Eksekusi):} Mobilisasi personel Polri dalam skala memadai untuk menjamin keamanan fisik jurusita dan mencegah benturan horizontal di lapangan;
    \item \textbf{Biaya Teknis Badan Pertanahan Nasional (BPN):} Biaya resmi pengukuran ulang, pematokan, dan pencocokan batas tanah (\textit{konstatering}) objek sengketa;
    \item \textbf{Biaya Jasa Penilai Publik (KJPP) dan Penilai Harta:} Honorarium penaksir harga independen berlisensi guna menilai harga wajar aset sita eksekusi sebelum proses pelelangan; serta
    \item \textbf{Biaya Operasional Lelang Eksekusi:} Administrasi dan publikasi risalah lelang melalui Kantor Pelayanan Kekayaan Negara dan Lelang (KPKNL).
\end{enumerate}
Ketiadaan pos anggaran negara untuk membiayai komponen-komponen eksternal tersebut mereduksi wibawa peradilan: putusan hakim yang diucapkan ``Demi Keadilan Berdasarkan Ketuhanan Yang Maha Esa'' terdegradasi menjadi sekadar dokumen hampa di atas kertas (\textit{illusory justice}) karena tidak berdaya dieksekusi secara nyata, mencederai hak konstitusional masyarakat atas kepastian hukum yang adil (Pasal 28D ayat (1) UUD NRI 1945).

Dalam kerangka perumusan haluan negara, usulan patokan eskalasi bertahap menuju kisaran 1,0\%--1,5\% APBN (berbasis postur belanja APBN 2025 sebesar Rp3.621,30 triliun hingga APBN 2026 sebesar Rp3.842,70 triliun dan RAPBN 2027 sebesar Rp4.106,26 triliun) harus dipahami sebagai \textbf{sebuah hipotesis kebijakan normatif dan benchmark arah haluan (PPHN)} jangka menengah-panjang. Sebagai catatan perkembangan fiskal, pasca-PP No. 42 Tahun 2025 pagu Mahkamah Agung pada TA 2026 memang meningkat ke kisaran Rp23 triliun, namun lonjakan tersebut terserap dominan bagi tunjangan jabatan hakim sehingga belum menyentuh kesetaraan tunjangan kepaniteraan maupun beban belanja modal infrastruktur secara utuh. Mengingat sempitnya ruang fiskal APBN akibat beban mandatori konstitusional, usulan ini tidak dimaksudkan sebagai \textit{mandatory spending} kaku yang dipaksakan seketika, melainkan sasaran bertahap (\textit{incremental phase-in}) berbasis \textit{costing} kebutuhan riil dan beban perkara. Pergeseran menuju kriteria pendanaan yang rasional, transparan, dan terukur berbasis beban perkara serta kebutuhan operasional satuan kerja terbukti krusial untuk menyeimbangkan otonomi fiskal dengan akuntabilitas kinerja layanan pengadilan tingkat pertama,\footnote{Federica Viapiana, ``A Performance-Based Budget in the Judiciary: Allocation of Resources and Performance Variability in First Instance Courts. An Analysis of Three Case Studies,'' \textit{Journal of Public Budgeting, Accounting \& Financial Management} 33, no. 2 (2021): 177--206, \url{https://doi.org/10.1108/JPBAFM-03-2020-0031}.} yang terdistribusi secara proporsional pada empat pos prioritas:
\begin{enumerate}
    \item \textbf{Belanja Pegawai dan Pemulihan Kesejahteraan Aparatur Peradilan (50\%--55\%):} Menutup disparitas ekstrem pasca-PP No. 42 Tahun 2025 melalui penyesuaian hak keuangan Panitera, Panitera Pengganti, Jurusita, dan administrator kesekretariatan secara berkeadilan;
    \item \textbf{Belanja Modal Infrastruktur dan Sarana Satker (25\%--30\%):} Menuntaskan rehabilitasi lebih dari 400 gedung pengadilan yang belum memenuhi standar prototipe (terutama pada satuan kerja daerah pemekaran dan wilayah 3T), rumah dinas, serta transportasi sidang keliling guna memutus ketergantungan hibah Pemda;
    \item \textbf{Modernisasi Peradilan Digital dan Operasional Teknis (10\%--12\%):} Penguatan keamanan siber, integrasi \textit{e-court}, dan digitalisasi pelacakan eksekusi (\textit{e-enforcement}); serta
    \item \textbf{Dana Talangan Operasional Eksekusi Perdata Berkeadilan / \textit{Civil Judgment Enforcement Fund} (3\%--5\%):} Pos khusus pembiayaan pengamanan Polri, konstatering BPN, penilai publik KJPP, dan publikasi lelang KPKNL bagi perkara prodeo serta tunggakan eksekusi kronis, dengan skema \textbf{dana bergulir (\textit{revolving fund})} bagi perkara komersial yang biayanya dipulihkan kembali ke kas negara dari hasil lelang aset tereksekusi; sementara untuk eksekusi riil atau pengosongan non-lelang bagi pemohon tidak mampu (\textit{prodeo}), dana ini bersifat belanja akses keadilan negara, sedangkan perkara komersial non-lelang tetap dibebankan kepada pihak tereksekusi melalui mekanisme pemulihan biaya (\textit{cost-shifting}).
\end{enumerate}

Dalam perkembangannya, narasi yang selama ini mengemuka bahwa hak keuangan hakim diabaikan selama lebih dari satu dekade memang telah direspons oleh pemerintah melalui penetapan Peraturan Pemerintah Nomor 42 Tahun 2025 yang ditindaklanjuti dengan Surat Edaran Direktur Jenderal Perbendaharaan Kementerian Keuangan Nomor SE-1/PB/2026.\footnote{Pemerintah RI, \textit{Peraturan Pemerintah Nomor 42 Tahun 2025 tentang Perubahan Keempat atas Peraturan Pemerintah Nomor 94 Tahun 2012 tentang Hak Keuangan dan Fasilitas Hakim yang Berada di Bawah Mahkamah Agung}, Lembaran Negara RI Tahun 2025 Nomor 154; Kementerian Keuangan RI, \textit{Surat Edaran Direktur Jenderal Perbendaharaan Nomor SE-1/PB/2026 tentang Perubahan Ketiga atas SE-38/PB/2012} (Jakarta: Ditjen Perbendaharaan, 2026).} Melalui regulasi eksekutif ini, tunjangan jabatan hakim dinaikkan secara drastis, mulai dari Rp46,7 juta bagi Hakim Pratama di Pengadilan Tingkat Pertama Kelas II hingga mencapai Rp110,5 juta bagi Ketua Pengadilan Tingkat Banding. Namun demikian, respons karikatif-parsial ini justru membongkar anomali dan distorsi struktural yang jauh lebih fundamental dalam relasi antarlembaga negara:

Pertama, \textbf{Disparitas Ekstrem Antar-Aparatur Peradilan}. Kenaikan tunjangan jabatan hakim tidak dibarengi perbaikan hak keuangan aparatur peradilan non-hakim (Panitera, Panitera Pengganti, Jurusita, dan staf). Cabang eksekutif memaknai norma ``Kekuasaan Kehakiman'' Pasal 24 UUD NRI 1945 secara reduksionis hanya berpusat pada personal ``Hakim'', bukan kesatuan sistem peradilan (\textit{the judiciary system}). Akibatnya, alokasi belanja aparatur terserap dominan bagi hakim, menyisakan porsi marjinal bagi aparatur lainnya. Padahal proses peradilan adalah orkestrasi teknis yudisial: putusan hakim mustahil dijatuhkan sah tanpa pencatatan berita acara akurat oleh Panitera, dan putusan mandul tanpa eksekusi lapangan oleh Jurusita yang menghadapi risiko fisik tinggi, di tengah ketimpangan beban kerja di satuan kerja besar—di mana satu Panitera Pengganti kerap merangkap mendampingi beberapa majelis hakim sekaligus tanpa insentif risiko persidangan, sedangkan Sekretaris selaku Kuasa Pengguna Anggaran memikul tanggung jawab pemeriksaan BPK atas pengelolaan DIPA dan BMN tanpa kompensasi risiko yang sepadan. Membiarkan panitera, jurusita, dan sekretariat terkurung dalam remunerasi ASN umum menciptakan disparitas destruktif dan melemahkan ketahanan integritas serta layanan peradilan.\footnote{Christina Clarissa Intania, op. cit., hlm. 4--10.} Lebih jauh lagi, ketimpangan kesejahteraan yang menganga ini memicu distorsi struktur sumber daya manusia berupa \textbf{pengurasan talenta internal (\textit{internal brain drain})}. Ketika pintu pengadaan calon hakim dibuka dari jalur ASN internal Mahkamah Agung, ribuan aparatur muda terbaik berbondong-bondong meninggalkan fungsi kepaniteraan semata demi mengejar lompatan status dan kesejahteraan. Akibatnya, struktur kepaniteraan terancam mengalami \textbf{seleksi merugikan (\textit{adverse selection})}—berubah dari profesi spesialis hukum acara menjadi sekadar ruang tunggu atau tempat bertumpuknya residu SDM. Padahal, kualitas putusan pengadilan tidak dapat bertumpu pada kecerdasan hakim semata, melainkan mensyaratkan kehadiran panitera berinteligensi hukum tinggi untuk menjamin ketelitian fakta persidangan, keabsahan berita acara, dan akurasi hukum acara.

Kedua, \textbf{Subordinasi Melalui Regulasi Eksekutif}. Seluruh penentuan hak keuangan peradilan masih dituangkan dalam instrumen hukum cabang eksekutif berupa Peraturan Pemerintah, Keputusan Presiden, hingga Surat Edaran Kementerian Keuangan. Bahkan penggunaan kembali Penerimaan Negara Bukan Pajak (PNBP) fungsional peradilan—seperti biaya proses perkara, leges, dan denda—tetap harus melewati persetujuan Kementerian Keuangan. Ketergantungan ini membuktikan bahwa cabang yudikatif masih diposisikan sebagai pemohon yang nasibnya ditentukan oleh kompromi eksekutif, bukan melalui kemandirian anggaran yang terlindungi secara konstitusional.

\subsection{Keterjebakan dalam Rezim UU ASN: Urgensi Pemisahan Korps Aparatur Sipil Peradilan (Judicial Civil Service)}

Kelemahan mendasar kedua terletak pada desain manajemen kepegawaian aparatur peradilan yang sarat distorsi doktriner:

\begin{enumerate}
    \item \textbf{Status Sui Generis Hakim Karier:} Berdasarkan Pasal 19 UU No. 48 Tahun 2009 serta Pasal 57 dan Pasal 58 UU No. 20 Tahun 2023 tentang ASN (yang mempertahankan substansi Pasal 121--122 UU No. 5 Tahun 2014) dan Putusan Mahkamah Konstitusi Nomor 32/PUU-XII/2014,\footnote{Mahkamah Konstitusi RI, \textit{Putusan Nomor 32/PUU-XII/2014}, perihal Pengujian UU ASN terhadap UUD 1945, diucapkan tanggal 20 April 2015, bagian Pertimbangan Hukum.} hakim ditegaskan sebagai Pejabat Negara yang terpisah dari PNS biasa. Namun, dalam tata kelola birokrasi peradilan, timbul anomali doktriner: hakim menyandang label ``pejabat negara'', tetapi administrasi teknis kepangkatan dan golongan ruangnya tetap dikelola BKN selayaknya birokrasi eksekutif. Karakteristik otentik Pejabat Negara dalam hukum tata negara lazimnya mengemban \textbf{masa jabatan periodik tertentu (\textit{fixed-term}, umumnya lima tahun)} serta memegang fungsi kepemimpinan publik atau mandat keterwakilan khusus. Secara doktrinal, karakter tersebut sejatinya lebih tepat melekat pada pimpinan lembaga peradilan—seperti Ketua dan Wakil Ketua Pengadilan, Ketua MA, dan Ketua MK—serta \textbf{Hakim Ad Hoc} bermasa jabatan lima tahunan, kendati dalam hukum positif pasca-Putusan MK Nomor 32/PUU-XII/2014 pengakuan status Pejabat Negara bagi Hakim Ad Hoc justru masih ditolak. Sebaliknya, hakim karier pada hakikatnya adalah \textbf{Pejabat Yudisial Negara Mandiri (\textit{Judicial Officers})} berkarakter birokrasi peradilan profesional tetap (\textit{tenured career judiciary}) hingga usia pensiun. Menyamakan hakim karier dengan pejabat negara politis tanpa undang-undang profesi independen terbukti hanya menjadi tameng administratif untuk melegitimasi kenaikan tunjangan parsial eksekutif. Hakim karier dan aparatur peradilan mutlak memerlukan kodifikasi profesi tersendiri (\textit{sui generis}) melalui pengesahan RUU Jabatan Hakim yang kini telah masuk Prolegnas Jangka Menengah 2025--2029 dan Prolegnas Prioritas 2026.
    \item \textbf{Ketimpangan Distribusi Beban Perkara, Krisis Majelis Kelas IA, dan Pola Residensi Yudisial:} Di berbagai negara bertradisi peradilan karier (\textit{career judiciary}) \textit{civil law}—seperti Jerman (\textit{Richter auf Probe} selama 3--5 tahun), Belanda (\textit{Rechterlijk Ambtenaar in Opleiding}), maupun Jepang (\textit{Assistant Judge}/\textit{Hanjiho} selama 10 tahun pertama yang duduk dalam majelis sebelum berwenang memutus mandiri)—penempatan awal hakim pemula dirancang berjenjang dalam komposisi majelis di pengadilan dengan variasi perkara memadai. Pola penempatan langsung hakim baru ke Pengadilan Kelas II di Indonesia selama ini menyisakan inefisiensi ganda. Kendati di sejumlah satuan kerja Kelas II hakim junior memang tidak selalu memutus sendiri karena masih dapat duduk sebagai anggota majelis bersama Ketua dan Wakil Ketua Pengadilan, volume dan kompleksitas perkara di Kelas II relatif terbatas; sebaliknya, satuan kerja \textbf{Pengadilan Kelas IA dan Kelas IA Khusus} yang menanggung beban ribuan perkara kompleks per tahun justru mengalami kekurangan formasi majelis hakim secara akut.

          Ketimpangan distribusi inilah yang memaksa meluasnya diskresi penerbitan \textbf{Surat Izin Hakim Tunggal} oleh Ketua Mahkamah Agung serta memicu beban kerja berlebih (\textit{overload}) bagi hakim senior di satker besar. Oleh karena itu, Indonesia mendesak mengadopsi skema \textbf{residensi yudisial (\textit{judicial residency})} selama 2--3 tahun dengan menempatkan hakim pemula terlebih dahulu sebagai \textbf{anggota majelis di Pengadilan Kelas IA dan Kelas IA Khusus}. Skema ini memberikan manfaat ganda: secara operasional langsung menambal defisit majelis hakim pada satker berbeban perkara tinggi guna menekan ketergantungan pada persidangan hakim tunggal, dan secara substantif mempercepat kematangan teknis hakim junior di bawah supervisi beragam hakim senior pada perkara-perkara rumit. Secara filosofis, keadilan bukanlah ajang coba-coba teori (\textit{trial and error}): kesalahan yudisial menyangkut hak asasi manusia tidak boleh ditoleransi (\textit{zero tolerance on judicial miscarriage}). Setelah menuntaskan masa residensi di Kelas IA/IA Khusus, barulah mereka didistribusikan secara berkeadilan ke Pengadilan Kelas II atau dipromosikan memimpin majelis, ditopang jaminan fasilitas rumah dinas dan insentif kemahalan yang layak.
    \item \textbf{Kebuntuan Nomenklatur KemenPAN-RB, Teknis Kepegawaian BKN, dan Dwitunggal Administrasi Peradilan:} Bukti konkret kegagalan penyeragaman peradilan ke dalam birokrasi eksekutif adalah fakta bahwa hingga kini \textbf{jabatan Panitera dan Jurusita tidak dapat dinomenklaturkan secara tepat oleh Kementerian PAN-RB} selaku otoritas regulasi jabatan fungsional ASN, sementara Badan Kepegawaian Negara (BKN) sebatas melaksanakan teknis pengelolaan kepegawaian (termasuk administrasi kepangkatan dan golongan ruang hakim yang berstatus pejabat negara). KemenPAN-RB mengalami kebuntuan struktural karena tata kelola peradilan bertumpu pada \textbf{dwitunggal kepemimpinan administrasi (\textit{twin pillars of court administration})}: Panitera memimpin administrasi teknis perkara yudisial, sedangkan Sekretaris memimpin kesekretariatan umum sebagai tulang penopang utama operasional, finansial, dan kelembagaan pengadilan yang mutlak wajib diperkokoh.\footnote{Mahkamah Agung RI, \textit{Peraturan Mahkamah Agung Nomor 7 Tahun 2015 tentang Organisasi dan Tata Kerja Kepaniteraan dan Kesekretariatan Peradilan}, Berita Negara Republik Indonesia Tahun 2015 Nomor 1532.} Jika Panitera dipaksakan masuk rumpun jabatan struktural ASN eksekutif, timbul kerancuan penentuan hierarki eselon karena Panitera bukan bawahan Sekretaris pengadilan, bahkan posisinya mendampingi majelis hakim di persidangan. Sebaliknya, jika diposisikan sebagai pejabat fungsional biasa, Panitera memegang wewenang komando eksekusi, memimpin jajaran kepaniteraan, dan menandatangani akta otentik demi hukum. Panitera dan Jurusita merupakan pejabat fungsional yudisial penegakan hukum (\textit{judicial auxiliary officers}) yang disumpah khusus, sehingga memaksakan mereka ke dalam rezim ASN sipil sama kelirunya dengan memaksakan prajurit militer atau aparat kepolisian ke dalam rezim ASN umum.

          Keterbelengguan rezim ASN ini turut mematikan tangga karier yudisial (\textit{judicial career ladder}). Panitera Pengganti yang setiap hari mengawal persidangan dan menguasai hukum acara teramputasi peluangnya beralih menjadi hakim. Kendati Pasal 5 ayat (1) huruf b Peraturan Mahkamah Agung Nomor 5 Tahun 2025 tentang Pengadaan Hakim Pengadilan Tingkat Pertama telah membuka pintu masuk calon hakim dari kalangan PNS Mahkamah Agung,\footnote{Mahkamah Agung RI, \textit{Peraturan Mahkamah Agung Nomor 5 Tahun 2025 tentang Pengadaan Hakim Pengadilan Tingkat Pertama}, diundangkan tanggal 31 Desember 2025, Pasal 5 ayat (1) huruf b.} peluang Panitera Pengganti senior tetap terganjal oleh batas kaku usia pengangkatan hakim pertama maksimal 40 tahun dalam Pasal 14 ayat (1) huruf h UU No. 49 Tahun 2009 (Peradilan Umum), Pasal 13 ayat (1) huruf i UU No. 50 Tahun 2009 (Peradilan Agama), dan Pasal 14 ayat (1) huruf g UU No. 51 Tahun 2009 (PTUN).\footnote{Ketiga undang-undang paket peradilan tersebut seragam mensyaratkan calon hakim karier berusia 25 hingga 40 tahun saat pengangkatan resmi di bawah MA.} Sejalan dengan Putusan Mahkamah Konstitusi Nomor 43/PUU-XIII/2015 yang menegaskan bahwa pengadaan hakim tingkat pertama merupakan kewenangan eksklusif Mahkamah Agung tanpa pelibatan lembaga eksternal dalam seleksi,\footnote{Mahkamah Konstitusi RI, \textit{Putusan Nomor 43/PUU-XIII/2015}, tanggal 7 Oktober 2015, yang membatalkan keterlibatan Komisi Yudisial dalam seleksi hakim tingkat pertama; lihat pula Muh. Ridha Hakim, ``Tafsir Independensi Kekuasaan Kehakiman dalam Putusan Mahkamah Konstitusi,'' \textit{Jurnal Hukum dan Peradilan} 7, no. 2 (2018): 279--296, \url{https://doi.org/10.25216/jhp.7.2.2018.279-296}.} skema alih jalur internal (\textit{in-service lateral entry}) ini mutlak dibuka melalui revisi tiga UU Peradilan dan dijalankan oleh panitia seleksi mandiri Mahkamah Agung berbasis meritokrasi, transparansi kepatuhan LHKPN, dan eksaminasi integritas perkara. Agar transisi karier berlangsung adil dan rasional secara organisasional, skema \textit{lateral entry} bagi Panitera Pengganti senior wajib disertai \textbf{pengakuan masa kerja (\textit{credit of service})}, penyetaraan golongan ruang dan kepangkatan yang telah dicapai, serta penyesuaian durasi pendidikan dan residensi awal mengingat penguasaan praktik hukum acara persidangan mereka telah teruji secara empiris. Kendati demikian, pembukaan jalur \textit{in-service lateral entry} tersebut wajib diimbangi dengan penataan remunerasi dan prestise profesi kepaniteraan melalui prinsip \textbf{keseimbangan jalur karier paralel (\textit{dual-track career parity})}. Tanpa kesetaraan insentif yang bermartabat bagi Panitera dan Jurusita, jalur alih profesi dari PNS Mahkamah Agung sebagaimana dibuka PERMA Nomor 5 Tahun 2025 justru berisiko berubah menjadi eksodus massal bermotif ekonomi yang menguras talenta unggul dari lini kepaniteraan. Dengan keseimbangan karier paralel, \textit{lateral entry} berfungsi murni sebagai saluran mobilitas meritokratis berbasis panggilan profesi, sementara aparatur berinteligensi tinggi tetap memiliki insentif rasional untuk berkarier penuh dan menjaga standar mutu teknis pada korps kepaniteraan.
\end{enumerate}

Sebagaimana negara telah secara tegas \textbf{memisahkan rezim manajemen kepegawaian ASN umum dengan Polri dan TNI} demi menjaga doktrin dan fungsi pertahanan-keamanan nasional, maka kekuasaan kehakiman mutlak menuntut pemisahan serupa. Berbeda dari TNI dan Polri yang tetap berada dalam cabang eksekutif di bawah Presiden, pemisahan korps Aparatur Sipil Peradilan justru jauh lebih mendesak secara konstitusional karena lembaga peradilan berdiri sebagai cabang kekuasaan mandiri yang setara dengan eksekutif. Gagasan ini sekaligus mengisi celah pembaruan dari \textit{Cetak Biru Pembaruan Peradilan 2010--2035} Mahkamah Agung RI dan \textit{Kajian Akademik Rekonseptualisasi Arah Penataan Kekuasaan Kehakiman} Badan Pengkajian MPR RI (2020) yang belum menuntaskan integrasi otonomi fiskal dan pemisahan korps aparatur secara simultan.\footnote{Mahkamah Agung RI, \textit{Cetak Biru Pembaruan Peradilan 2010--2035} (Jakarta: Mahkamah Agung RI, 2010), hlm. 35--42; Badan Pengkajian MPR RI, \textit{Kajian Akademik Rekonseptualisasi Arah Penataan Kekuasaan Kehakiman} (Jakarta: Sekretariat Jenderal MPR RI, 2020), hlm. 88--94.} Pelembagaan korps \textbf{Aparatur Sipil Peradilan (\textit{Judicial Civil Service})} yang mandiri di bawah kewenangan eksklusif Mahkamah Agung—dengan pengawasan sistem merit independen—merupakan keniscayaan konstitusional guna memutus loyalitas ganda birokrasi terhadap kementerian eksekutif (KemenPAN-RB/BKN) dan mengamankan integritas keadilan. Adapun bagi lingkungan Peradilan Militer, personel hakim dan kepaniteraan militernya tetap tunduk pada kekhususan pembinaan keprajuritan TNI secara fungsional, sementara aparatur sipil kesekretariatannya terintegrasi ke dalam korps \textit{Judicial Civil Service}.

Dalam konteks relasi anggaran, pembenahan ini mentransformasikan relasi subordinatif yudikatif-eksekutif menjadi negosiasi konstitusional langsung antara Mahkamah Agung dan DPR RI selaku pemegang hak bujet (Pasal 20A ayat (1) UUD NRI 1945). Melalui mekanisme \textit{direct budget submission} ala Amerika Serikat, rencana anggaran dibahas secara transparan oleh pejabat administrator kesekretariatan pengadilan (\textit{court managers}) dalam forum dengar pendapat resmi parlemen, bukan melalui lobi tertutup hakim pemutus perkara. Ditopang alokasi bertahap menuju 1,0\%--1,5\% APBN (\textit{statutory floor} ala Kosta Rika) guna membiayai perbaikan lebih dari 400 gedung pengadilan yang belum memenuhi standar prototipe (terutama pada satuan kerja daerah pemekaran dan wilayah 3T), hak keuangan berkeadilan aparatur, penuntasan tunggakan eksekusi perdata nasional (Polri, BPN, KJPP, KPKNL), serta garansi \textit{non-diminution rule} ala Filipina, kemandirian yudikatif terlindungi utuh dari pemangkasan eksekutif maupun jebakan transaksi politik.

\subsection{Kajian Komparatif Internasional: Model Penjaminan Otonomi Anggaran dan Aparatur Yudikatif}

Untuk merumuskan desain hukum yang kokoh, penting untuk melakukan perbandingan dengan negara-negara yang telah berhasil melindungi kemerdekaan peradilannya melalui garansi konstitusional dan legislasi khusus:\footnote{Lihat kajian perbandingan penganggaran dan indikator kemandirian peradilan dalam Federica Viapiana, op. cit., hlm. 177--206; serta Frans van Dijk dan Geoffrey Vos, op. cit., hlm. 1--21.}

\subsubsection{Kosta Rika: Model Garansi Alokasi Anggaran Minimum Konstitusional}
Konstitusi Kosta Rika Tahun 1949 (\textit{Constitución Política de la República de Costa Rica}) memberikan model perlindungan fiskal paling tegas di dunia. Pasal 177 Konstitusi Kosta Rika secara eksplisit mewajibkan negara mengalokasikan anggaran tahunan untuk kekuasaan kehakiman (\textit{Poder Judicial}) sekurang-kurangnya 6\% dari total pendapatan biasa negara (\textit{ordinary revenues}):
\begin{quote}
    ``\textit{...an amount not less than six percent of the ordinary income calculated for the fiscal year shall be earmarked for the Judicial Power...}''\footnote{Constitución Política de la República de Costa Rica (1949), Article 177.}
\end{quote}
Jaminan \textit{mandatory spending} berbasis konstitusi ini telah membentengi cabang kehakiman Kosta Rika dari pembalasan politik eksekutif dan parlemen. Namun secara perbandingan hukum, perlu dicatat bahwa \textit{Poder Judicial} Kosta Rika membiayai bukan hanya badan peradilan, melainkan mencakup Kejaksaan (\textit{Ministerio P\'{u}blico}), Kepolisian Investigasi Yudisial (\textit{OIJ}), dan Lembaga Bantuan Hukum Publik (\textit{Defensa P\'{u}blica}). Oleh karena itu, patokan hipotetis bertahap menuju kisaran 1,0\%--1,5\% APBN bagi Mahkamah Agung RI merupakan proporsi yang sangat terukur dan rasional, mengingat kepolisian dan kejaksaan di Indonesia telah memiliki pagu tersendiri di luar MA.

\subsubsection{Filipina: Otonomi Fiskal dan Klausul Non-Diminution Anggaran}
Filipina memberikan preseden penting di kawasan Asia Tenggara melalui Konstitusi 1987 (\textit{1987 Constitution of the Republic of the Philippines}). \textit{Article VIII, Section 3} menetapkan jaminan otonomi fiskal peradilan dengan klausul larangan penurunan anggaran:
\begin{quote}
    ``\textit{The Judiciary shall enjoy fiscal autonomy. Appropriations for the Judiciary may not be reduced by the legislature below the amount appropriated for the previous year and, after approval thereof, shall be automatically and regularly released.}''\footnote{The 1987 Constitution of the Republic of the Philippines, Article VIII, Section 3.}
\end{quote}
Pasal ini mengandung tiga garansi konstitusional sekaligus:
\begin{enumerate}
    \item Pengakuan formal atas otonomi fiskal kekuasaan kehakiman (\textit{enjoy fiscal autonomy});
    \item Larangan bagi lembaga legislatif untuk memotong anggaran peradilan lebih rendah daripada alokasi tahun sebelumnya (\textit{non-diminution rule}); dan
    \item Kewajiban bagi bendahara negara eksekutif untuk mencairkan anggaran yang telah disahkan secara otomatis dan berkala (\textit{automatic and regular release}) tanpa syarat persetujuan DIPA kementerian keuangan, dengan tetap tunduk pada audit eksternal perbendaharaan.
\end{enumerate}

\subsubsection{Amerika Serikat: Mekanisme Pengajuan Anggaran Langsung (Direct Budget Submission)}
Di tingkat federal Amerika Serikat, perlindungan anggaran yudikatif diatur melalui mekanisme perundang-undangan dalam \textit{Budget and Accounting Act 1921} (dikodifikasikan dalam 31 U.S.C. \S\ 1105(b)). Berdasarkan aturan ini, cabang eksekutif melalui \textit{Office of Management and Budget} (OMB) dilarang mengubah, mengoreksi, atau memotong usulan anggaran yang diajukan oleh cabang yudikatif. Usulan anggaran dirumuskan secara mandiri oleh \textit{Judicial Conference of the United States} (dipimpin oleh Ketua Mahkamah Agung AS) dan diadministrasikan oleh \textit{Administrative Office of the United States Courts} (AO). Presiden diwajibkan oleh undang-undang untuk meneruskan draf anggaran yudikatif tersebut secara utuh tanpa modifikasi apa pun langsung ke Kongres Amerika Serikat (\textit{direct budget submission}).\footnote{Congressional Research Service, \textit{The President's Budget: Overview of Structure and Timing of Submission to Congress}, CRS Report R43163 (Washington, D.C.: Library of Congress, 25 Juli 2013); lihat pula 31 U.S.C. \S\ 1105(b).} Selain itu, staf peradilan federal di AS berada dalam kategori \textit{excepted service} di luar kendali badan kepegawaian eksekutif (\textit{Office of Personnel Management}), melainkan diatur langsung oleh sistem peradilan (\textit{Judicial Branch Personnel System}).

% =========================================================================
% 4. REKOMENDASI OPERASIONAL BAGI PPHN
% =========================================================================
\section{Rekomendasi Operasional bagi Pokok-Pokok Haluan Negara (PPHN)}

\begin{table*}[t]
    \centering
    \caption{Matriks Rekomendasi Operasional bagi PPHN dan Tahapan Implementasi Kebijakan}
    \label{tab:rekomendasi_pphn}
    \footnotesize
    \renewcommand{\arraystretch}{1.12}
    \begin{tabularx}{\textwidth}{@{} >{\raggedright\arraybackslash}p{2.4cm} >{\raggedright\arraybackslash}X >{\raggedright\arraybackslash}p{4.6cm} >{\centering\arraybackslash}p{2.4cm} @{}}
        \toprule
        \textbf{Dimensi}                      & \textbf{Materi Muatan Haluan (PPHN)}                                                                                                                        & \textbf{Instrumen Kebijakan \& Tindak Lanjut Regulasi}                                                                                                                                     & \textbf{Jangka Waktu Implementasi}                  \\
        \midrule
        \textbf{Otonomi Fiskal}               & Menggariskan jaminan kemandirian fiskal melalui mekanisme penerusan rancangan anggaran peradilan secara utuh ke DPR RI tanpa modifikasi eksekutif (\textit{transmission without modification} selaras Pasal 23 ayat (2) UUD NRI 1945) serta larangan pemotongan sepihak. & $\bullet$~Revisi UU No. 17/2003 tentang Keuangan Negara \par\vspace{2pt} $\bullet$~Revisi UU Mahkamah Agung & \textbf{Jangka Pendek}\par (2026--2027) \\
        \addlinespace[4pt]
        \textbf{Alokasi Minimum}              & Menetapkan mandat alokasi anggaran tetap kekuasaan kehakiman secara bertahap menuju kisaran 1,0\%--1,5\% APBN, pelembagaan Dana Penuntasan Tunggakan Eksekusi Perdata Nasional, dan klausul \textit{non-diminution}. & $\bullet$~Pembentukan klausul norma dalam RUU APBN \par\vspace{2pt} $\bullet$~Pengesahan RUU Otonomi Peradilan \& Regulasi Eksekusi Perdata Berkeadilan & \textbf{Jangka Menengah}\par (2027--2029)           \\
        \addlinespace[4pt]
        \textbf{Sistem \& Aparatur Peradilan} & Menginstruksikan reorientasi norma organisasi sistem pengadilan serta pelembagaan korps Aparatur Sipil Peradilan (\textit{Judicial Civil Service}) mandiri. & $\bullet$~Revisi UU No. 48/2009 tentang Kekuasaan Kehakiman \par\vspace{2pt} $\bullet$~Pengesahan RUU Jabatan Hakim \par\vspace{2pt} $\bullet$~Revisi UU ASN \& Harmonisasi 3 UU Peradilan & \textbf{Jangka Pendek \& Menengah}\par (2026--2029) \\
        \addlinespace[4pt]
        \textbf{Desain Konstitusi}            & Menata ulang relasi horizontal trias politika dan memasukkan garansi otonomi peradilan ke dalam materi konstitusi.                                          & $\bullet$~Rekomendasi Badan Pengkajian MPR RI untuk Agenda Perubahan Kelima UUD NRI 1945 (Pasal 24 ayat (4))                                                                               & \textbf{Jangka Panjang}\par (2029--2034)            \\
        \bottomrule
    \end{tabularx}
    \vspace{4pt}
    \begin{minipage}{\textwidth}
        \footnotesize \textit{Sumber: Formulasi Konseptual Penulis untuk Bahan Pertimbangan Badan Pengkajian MPR RI (2026).}
    \end{minipage}
\end{table*}

Penataan relasi horizontal antara cabang yudikatif, eksekutif, dan legislatif bukan sekadar isu perbaikan manajerial internal Mahkamah Agung, melainkan agenda fundamental ketatanegaraan. Prinsip otonomi fiskal dan kemandirian aparatur ini relevan diperluas secara proporsional bagi kesatuan cabang kekuasaan kehakiman dalam Bab IX UUD NRI Tahun 1945 (termasuk Mahkamah Konstitusi serta Komisi Yudisial dalam kapasitas kelembagaannya). Majelis Permusyawaratan Rakyat (MPR) melalui mandat pengkajian sistem ketatanegaraan dan penyusunan Pokok-Pokok Haluan Negara (PPHN) memegang peran vital mengintegrasikan otonomi institusional peradilan ke dalam materi haluan negara demi kesinambungan pembangunan berlandaskan supremasi konstitusi.\footnote{Lihat Badan Pengkajian MPR RI, \textit{Kajian Rancangan Substansi Pokok-Pokok Haluan Negara (PPHN)} (Jakarta: Sekretariat Jenderal MPR RI, 2025); serta Badan Pengkajian MPR RI, \textit{Kajian Bentuk Hukum Pokok-Pokok Haluan Negara (PPHN)} (Jakarta: Sekretariat Jenderal MPR RI, 2026).} Matriks rekomendasi operasional bagi PPHN dan agenda tindak lanjut legislasi dipaparkan secara komprehensif pada Tabel \ref{tab:rekomendasi_pphn}.

Efektivitas haluan otonomi peradilan tersebut sangat ditentukan oleh pilihan bentuk hukum PPHN yang ditetapkan MPR RI. Pasca-Putusan Mahkamah Konstitusi Nomor 66/PUU-XXI/2023 yang menegaskan bahwa MPR tidak lagi memiliki kewenangan membentuk Ketetapan MPR bersifat mengatur (\textit{regeling}) di luar amanat eksplisit UUD NRI Tahun 1945,\footnote{Muchamad Ali Safa'at, ``Akhir Perdebatan Kewenangan MPR Membentuk Ketetapan yang Bersifat Regeling (Studi Putusan Mahkamah Konstitusi Nomor 66/PUU-XXI/2023),'' \textit{Jurnal Konstitusi} 22, no. 4 (2025): 777--797, \url{https://doi.org/10.31078/jk2248}.} pengintegrasian prinsip pengarah kemandirian fiskal dan aparatur yudisial idealnya didudukkan melalui penguatan landasan konstitusional PPHN (\textit{Directive Principles of State Policies}) atau melalui Ketetapan MPR yang memiliki dasar konstitusional eksplisit agar memiliki daya ikat direktif terhadap DPR RI dan Presiden dalam penyusunan APBN maupun Prolegnas, bukan sekadar imbauan politik tanpa daya paksa ketatanegaraan.

Berikut adalah rumusan materi muatan dan rekomendasi operasional bagi PPHN pada Bidang Politik, Hukum, Pertahanan, dan Keamanan:

\subsection{Haluan Penataan Otonomi Anggaran Lembaga Yudikatif}
PPHN perlu memuat arahan direktif yang menetapkan arah kebijakan reformasi anggaran kekuasaan kehakiman sebagai berikut:
\begin{enumerate}
    \item \textbf{Alokasi Anggaran Minimum} (\textit{Statutory Floor}): Menetapkan haluan negara yang mewajibkan jaminan undang-undang alokasi anggaran tetap kekuasaan kehakiman secara bertahap menuju kisaran 1,0\%--1,5\% dari belanja negara APBN guna mendanai belanja modal lebih dari 400 gedung pengadilan yang belum memenuhi standar prototipe (terutama pada satuan kerja daerah pemekaran dan wilayah 3T) serta hak keuangan aparatur.
    \item \textbf{Pengajuan Anggaran Tanpa Pemotongan Eksekutif (\textit{Direct Budget Submission / Transmission Without Modification}):} Menetapkan mekanisme penyampaian rancangan anggaran kekuasaan kehakiman yang harmonis dengan Pasal 23 ayat (2) UUD NRI Tahun 1945. Pada tahap jangka pendek melalui revisi UU Keuangan Negara dan UU Mahkamah Agung (sebelum amandemen konstitusi), Presiden tetap menjadi penyampai formal RUU APBN ke DPR RI, tetapi diwajibkan meneruskan rancangan anggaran Mahkamah Agung secara utuh tanpa pemotongan, penyaringan pagu indikatif, maupun modifikasi sepihak oleh Kementerian Keuangan dan Bappenas (\textit{transmission without modification}, analog dengan 31 U.S.C. \S\ 1105(b)), termasuk fleksibilitas pemanfaatan PNBP fungsional peradilan secara langsung bagi kebutuhan teknis persidangan. Selanjutnya, rancangan tersebut dibahas secara terbuka antara pejabat administrator kesekretariatan peradilan (\textit{court managers}) dan DPR RI guna mengakhiri praktik lobi anggaran tertutup oleh hakim pemutus perkara, sebelum ditransformasikan menjadi kewenangan pengajuan langsung secara konstitusional pada agenda Perubahan Kelima UUD NRI Tahun 1945.
    \item \textbf{Perlindungan Berkelanjutan} (\textit{Non-Diminution Rule}): Melarang persetujuan anggaran peradilan lebih rendah dari tahun sebelumnya, serta mengecualikan kekuasaan kehakiman dari pemotongan anggaran sepihak (\textit{automatic adjustment}).
    \item \textbf{Alokasi Khusus Penuntasan Tunggakan Eksekusi Perdata Nasional (\textit{Civil Judgment Enforcement Fund}):} Menetapkan arahan penyediaan pos belanja operasional negara secara terukur untuk menuntaskan ribuan krisis tunggakan eksekusi perdata di seluruh lingkungan peradilan—meliputi biaya bantuan pengamanan Kepolisian RI (PAM Eksekusi), biaya teknis pengukuran dan \textit{konstatering} BPN, jasa penilai harga publik independen (KJPP), serta operasional lelang KPKNL. Guna menjaga disiplin fiskal APBN dan mencegah bahaya moral (\textit{moral hazard}), skema pembiayaan dirancang berdiferensiasi: pembiayaan penuh oleh negara dikhususkan bagi masyarakat pencari keadilan tidak mampu (\textit{prodeo}) dan tunggakan eksekusi kronis demi menegakkan kepastian hukum hakiki (termasuk eksekusi riil non-lelang), sedangkan untuk perkara perdata komersial dijalankan melalui mekanisme \textbf{dana bergulir (\textit{revolving fund})} yang seluruh biayanya ditarik dan dipulihkan kembali ke kas negara dari hasil pelelangan aset tereksekusi atau melalui pembebanan biaya (\textit{cost-shifting}) kepada pihak tereksekusi.
\end{enumerate}

\subsection{Klausul Haluan Pemisahan Rezim Aparatur dan Penataan Sistem Peradilan}
PPHN harus menggariskan peta jalan transformasi tata kelola aparatur dan kelembagaan peradilan, mencakup:
\begin{enumerate}
    \item \textbf{Reorientasi Paradigma UU Kekuasaan Kehakiman:} Menginstruksikan revisi Undang-Undang Nomor 48 Tahun 2009 tentang Kekuasaan Kehakiman agar muatannya tidak lagi bercorak hakim-sentris (\textit{judge-centric}), melainkan bertransformasi memperkuat keutuhan \textbf{Sistem dan Organisasi Pengadilan (\textit{court system and judicial administration})}. Regulasi ini harus meletakkan kemerdekaan peradilan sebagai satu kesatuan ekosistem institusional yang mencakup otonomi anggaran peradilan, perlindungan sarana persidangan, serta pengakuan kedudukan fungsional aparatur teknis yudisial lainnya (Panitera dan Jurusita) sebagai pilar integritas putusan.
    \item \textbf{Pengesahan RUU Jabatan Hakim dan Harmonisasi Tiga UU Peradilan:} Menginstruksikan pembentukan UU Jabatan Hakim yang mendudukkan hakim karier sebagai Pejabat Yudisial Negara Mandiri (\textit{sui generis Judicial Officers}), pola pembinaan bertingkat melalui residensi yudisial 2--3 tahun sebagai hakim anggota di Pengadilan Kelas IA dan Kelas IA Khusus guna mengatasi defisit majelis hakim berbeban perkara tinggi sekaligus mematangkan kompetensi teknis sebelum memimpin persidangan mandiri, serta harmonisasi batas usia pengangkatan hakim dalam UU Peradilan Umum, UU Peradilan Agama, dan UU PTUN dengan membuka skema alih promosi internal (\textit{in-service lateral entry}) bagi Panitera Pengganti berpengalaman—lengkap dengan jaminan \textbf{pengakuan masa kerja (\textit{credit of service})}, penyetaraan golongan ruang, dan penyesuaian masa pembinaan awal—melalui panel seleksi independen Mahkamah Agung berbasis meritokrasi, transparansi kepatuhan LHKPN, dan rekam jejak integritas etik.
    \item \textbf{Pembentukan Korps Aparatur Sipil Peradilan (\textit{Judicial Civil Service}):} Mengakhiri kebuntuan nomenklatur jabatan fungsional di KemenPAN-RB serta ketergantungan teknis administrasi kepegawaian (termasuk kepangkatan golongan ruang hakim) di BKN, dan menyelaraskan arsitektur dwitunggal administrasi pengadilan (\textit{twin pillars of court administration}) melalui pemisahan korps aparatur peradilan mandiri di bawah kewenangan eksklusif Mahkamah Agung (analog dengan pemisahan rezim TNI dan Polri). Desain korps mandiri ini memastikan Panitera dan Sekretaris memiliki pola jenjang karier struktural yang sama-sama berjalan seiring dengan kebutuhan lembaga, memperkokoh peran mutlak kesekretariatan sebagai tulang penopang operasional pengadilan, sekaligus memuliakan fungsi kepaniteraan teknis persidangan. Dengan struktur karier mandiri dan remunerasi yang kompetitif tersebut, jabatan kepaniteraan tidak lagi terdegradasi menjadi residu dari seleksi calon hakim, melainkan tumbuh sebagai jalur profesi hukum bergengsi yang mampu merekrut sekaligus mempertahankan SDM berinteligensi tinggi demi menopang mutu putusan peradilan.

          Pasca-pemisahan tersebut, hubungan antara Biro Kepegawaian (atau Badan Kepegawaian Yudikatif) Mahkamah Agung, KemenPAN-RB, dan BKN ditata ulang dari relasi subordinasi eksekutif menjadi \textbf{koordinasi antarcabang kekuasaan dan interoperabilitas data nasional}. Biro Kepegawaian MA memegang penuh otoritas penetapan nomenklatur jabatan teknis peradilan, rekrutmen, mutasi, serta kepangkatan dan golongan ruang tanpa memerlukan pertimbangan teknis BKN; KemenPAN-RB hanya berperan sebagai mitra koordinasi kuota formasi makro terhadap kapasitas APBN; sedangkan BKN terbatas pada fungsi interoperabilitas basis data kepegawaian negara (pencegahan gaji ganda dan sinkronisasi data pensiun Taspen/Kemenkeu), dengan tetap tunduk pada pengawasan eksternal lembaga pengawas sistem merit independen sesuai mandat Putusan MK Nomor 121/PUU-XXII/2024.
    \item \textbf{Skema Indeksasi Otomatis Hak Finansial:} Menetapkan formula penyesuaian berkala hak keuangan seluruh ekosistem peradilan secara otomatis berbasis inflasi guna mencegah disparitas dan distorsi kebijakan parsial eksekutif.
\end{enumerate}

\subsection{Formulasi Klausul Siap Adopsi bagi Rancangan Pokok-Pokok Haluan Negara (PPHN)}

Guna mempermudah Badan Pengkajian MPR RI mengintegrasikan rekomendasi ini ke dalam naskah haluan pembangunan hukum nasional, dirumuskan klausul direktif siap adopsi sebagai berikut:

\begin{quote}
\small
\textit{``Memperkokoh kemerdekaan kekuasaan kehakiman sebagai pilar negara hukum yang berkeadilan dan akuntabel melalui: (a) pelembagaan otonomi fiskal kekuasaan kehakiman secara bertahap dan terukur berbasis kebutuhan riil beban peradilan, mekanisme pengajuan anggaran langsung (direct budget submission) kepada DPR RI, perlindungan dari pemotongan anggaran sepihak (non-diminution rule), jaminan pembiayaan operasional negara atas efektivitas pelaksanaan putusan pengadilan yang berkekuatan hukum tetap, serta transparansi audit keuangan publik; dan (b) penataan sistem organisasi peradilan melalui pembentukan korps Aparatur Sipil Peradilan (judicial civil service) yang mandiri di bawah Mahkamah Agung dengan pengawasan sistem merit independen, kodifikasi profesi pejabat yudisial, pembinaan hakim berjenjang, serta jaminan tangga karier dan kesejahteraan yang berkeadilan bagi seluruh ekosistem aparatur peradilan.''}
\end{quote}

% =========================================================================
% 5. KESIMPULAN
% =========================================================================
\section{Kesimpulan}

Kajian ini membuktikan bahwa arsitektur ketatanegaraan pascareformasi menyisakan anomali fundamental dalam relasi horizontal trias politika. Berdasarkan konsistensi rumusan masalah dan hasil telaah, ditarik tiga butir kesimpulan:

\begin{enumerate}[topsep=2pt, itemsep=1.5pt, parsep=0pt]
    \item Kebijakan Satu Atap (\textit{one-roof system}) belum mewujudkan independensi peradilan sejati akibat subordinasi fiskal di bawah rezim UU Keuangan Negara dan penundukan aparatur non-hakim di bawah rezim UU ASN. Respons eksekutif melalui PP No. 42/2025 yang menaikkan tunjangan hakim sepihak justru memicu disparitas kesejahteraan internal yang ekstrem, kebuntuan nomenklatur jabatan fungsional kepaniteraan/kejurusitaan di KemenPAN-RB, serta ketergantungan teknis kepangkatan golongan ruang hakim maupun aparatur peradilan pada BKN.
    \item Preseden komparatif Kosta Rika, Filipina, dan Amerika Serikat membuktikan bahwa independensi yudikatif mutlak bertumpu pada otonomi anggaran (alokasi bertahap menuju kisaran 1,0\%--1,5\% APBN, pengajuan mandiri \textit{direct submission} ke parlemen, klausul \textit{non-diminution}, serta penyediaan dana penuntasan tunggakan eksekusi perdata nasional) dan pemisahan korps Aparatur Sipil Peradilan (\textit{Judicial Civil Service}) di bawah kendali eksklusif Mahkamah Agung—analog dengan pemisahan rezim TNI dan Polri dari ASN umum.
    \item Badan Pengkajian MPR RI perlu memosisikan agenda ini ke dalam Pokok-Pokok Haluan Negara (PPHN) sebagai haluan pembaruan hukum nasional, yang ditindaklanjuti bertahap melalui revisi UU Kekuasaan Kehakiman (berorientasi sistem pengadilan), revisi tiga paket UU Peradilan tahun 2009, pengesahan RUU Jabatan Hakim, revisi UU Keuangan Negara, revisi UU ASN, dan penguatan norma Pasal 24 UUD NRI Tahun 1945.
\end{enumerate}

% =========================================================================
% DAFTAR PUSTAKA
% =========================================================================
\section*{Daftar Pustaka}
\addcontentsline{toc}{section}{Daftar Pustaka}
\singlespacing
\fontsize{7.5}{8.8}\selectfont
\setlength{\parskip}{0pt}

\vspace{0.2em}
\noindent \textbf{Buku dan Bab Buku}
\begin{hangparas}{0.4cm}{1}
    Asshiddiqie, Jimly. 2021. \textit{Konstitusi dan Konstitusionalisme Indonesia}. Edisi Kedua. Jakarta: Sinar Grafika.

    Dixon, Rosalind, dan David Landau. 2021. \textit{Abusive Constitutional Borrowing: Legal Globalization and the Subversion of Liberal Democracy}. Oxford: Oxford University Press.

    Hamilton, Alexander, James Madison, dan John Jay. 1961. \textit{The Federalist Papers}. Diedit oleh Clinton Rossiter. New York: New American Library.

    Isra, Saldi. 2020. \textit{Pergeseran Kekuasaan Legislatif: Mengurai Hubungan Presiden dan DPR Pasca Perubahan UUD 1945}. Jakarta: Rajawali Pers.

    Levy, Jacob T. 2025. ``The Separation of Powers.'' Dalam Richard Bellamy dan Jeff King (eds.), \textit{The Cambridge Handbook of Constitutional Theory}, hlm. 279--296. Cambridge: Cambridge University Press.

    Manan, Bagir. 2009. \textit{Menegakkan Hukum Suatu Pencarian}. Jakarta: Asosiasi Advokat Indonesia.

    Marzuki, Peter Mahmud. 2017. \textit{Penelitian Hukum}. Edisi Revisi. Jakarta: Kencana Prenada Media Group.

    Mochtar, Zainal Arifin. 2021. \textit{Lembaga Negara Independen: Dinamika Perkembangan dan Urgensinya dalam Struktur Ketatanegaraan Indonesia}. Depok: Rajawali Pers.

    Montesquieu. 1900. \textit{The Spirit of Laws}. Diterjemahkan oleh Thomas Nugent. New York: Colonial Press.

    Pompe, Sebastiaan. 2005. \textit{The Indonesian Supreme Court: A Study of Institutional Collapse and Reconstruction}. Ithaca: Southeast Asia Program Publications, Cornell University.

    Shetreet, Shimon, dan Sophie Turenne. 2013. \textit{Judges on Trial: A Study of the Appointment and Accountability of the English Judiciary}. Edisi Kedua. Cambridge: Cambridge University Press.
\end{hangparas}

\vspace{0.15em}
\noindent \textbf{Jurnal}
\begin{hangparas}{0.4cm}{1}
    Bachmid, Fahri, dan Diani Indah Rachmitasari. 2026. ``Rekonstruksi Kemandirian Anggaran Kekuasaan Kehakiman dalam Sistem Ketatanegaraan Indonesia.'' \textit{Jurnal Legislasi Indonesia} 23, no. 2: 263--278. \url{https://doi.org/10.54629/jli.v23i2.1733}.

    Douglas, James W., dan Roger E. Hartley. 2003. ``The Politics of Court Budgeting in the States: Is Judicial Independence Threatened by the Budgetary Process?'' \textit{Public Administration Review} 63, no. 4: 441--454. \url{https://doi.org/10.1111/1540-6210.00307}.

    Hakim, Muh. Ridha. 2018. ``Tafsir Independensi Kekuasaan Kehakiman dalam Putusan Mahkamah Konstitusi.'' \textit{Jurnal Hukum dan Peradilan} 7, no. 2: 279--296. \url{https://doi.org/10.25216/jhp.7.2.2018.279-296}.

    Permadi, Restu, dan Fifiana Wisnaeni. 2020. ``Tinjauan Hukum Kemandirian dan Independensi Mahkamah Agung di dalam Sistem Ketatanegaraan Indonesia.'' \textit{Jurnal Pembangunan Hukum Indonesia} 2, no. 3: 399--415. \url{https://doi.org/10.14710/jphi.v2i3.399-415}.

    Prasetianingsih, Rahayu. 2011. ``Akuntabilitas Kekuasaan Kehakiman.'' \textit{Jurnal Konstitusi} 8, no. 5: 829--848. \url{https://doi.org/10.31078/jk858}.

    Safa'at, Muchamad Ali. 2025. ``Akhir Perdebatan Kewenangan MPR Membentuk Ketetapan yang Bersifat Regeling (Studi Putusan Mahkamah Konstitusi Nomor 66/PUU-XXI/2023).'' \textit{Jurnal Konstitusi} 22, no. 4: 777--797. \url{https://doi.org/10.31078/jk2248}.

    van Dijk, Frans, dan Geoffrey Vos. 2018. ``A Method for Assessment of the Independence and Accountability of the Judiciary.'' \textit{International Journal for Court Administration} 9, no. 3: 1--21. \url{https://doi.org/10.18352/ijca.276}.

    Viapiana, Federica. 2021. ``A Performance-Based Budget in the Judiciary: Allocation of Resources and Performance Variability in First Instance Courts. An Analysis of Three Case Studies.'' \textit{Journal of Public Budgeting, Accounting \& Financial Management} 33, no. 2: 177--206. \url{https://doi.org/10.1108/JPBAFM-03-2020-0031}.
\end{hangparas}

\vspace{0.15em}
\noindent \textbf{Laporan dan Dokumen Resmi}
\begin{hangparas}{0.4cm}{1}
    Badan Pengkajian MPR RI. 2020. \textit{Kajian Akademik Rekonseptualisasi Arah Penataan Kekuasaan Kehakiman}. Jakarta: Sekretariat Jenderal MPR RI.

    Badan Pengkajian MPR RI. 2025. \textit{Kajian Rancangan Substansi Pokok-Pokok Haluan Negara (PPHN)}. Jakarta: Sekretariat Jenderal MPR RI.

    Badan Pengkajian MPR RI. 2026. \textit{Kajian Bentuk Hukum Pokok-Pokok Haluan Negara (PPHN)}. Jakarta: Sekretariat Jenderal MPR RI.

    Congressional Research Service. 2013. \textit{The President's Budget: Overview of Structure and Timing of Submission to Congress}. CRS Report R43163. Washington, D.C.: Library of Congress.

    Intania, Christina Clarissa. 2025. \textit{Policy Assessment: Dampak Kebijakan Efisiensi Anggaran terhadap Independensi Peradilan}. Jakarta: The Indonesian Institute.

    Kementerian Keuangan RI. 2022--2025. \textit{Informasi APBN Tahun Anggaran 2022--2025}. Jakarta: Direktorat Jenderal Anggaran Kemenkeu RI.

    Mahkamah Agung RI. 2010. \textit{Cetak Biru Pembaruan Peradilan 2010--2035}. Jakarta: Mahkamah Agung RI.

    Mahkamah Agung RI. 2021. \textit{Laporan Tahunan Mahkamah Agung RI Tahun 2020}. Jakarta: Mahkamah Agung RI.

    Mahkamah Agung RI. 2022. \textit{Laporan Tahunan Mahkamah Agung RI Tahun 2021}. Jakarta: Mahkamah Agung RI.

    Mahkamah Agung RI. 2023. \textit{Laporan Tahunan Mahkamah Agung RI Tahun 2022}. Jakarta: Mahkamah Agung RI.

    Mahkamah Agung RI. 2024. \textit{Laporan Tahunan Mahkamah Agung RI Tahun 2023: Integritas Kuat, Peradilan Bermartabat}. Jakarta: Mahkamah Agung RI.

    Mahkamah Agung RI. 2025. \textit{Laporan Tahunan Mahkamah Agung RI Tahun 2024}. Jakarta: Mahkamah Agung RI.

    Mahkamah Agung RI. 2026a. \textit{Laporan Keuangan Mahkamah Agung RI Audited Tahun Anggaran 2025}. Jakarta: Mahkamah Agung RI.

    Mahkamah Agung RI. 2026b. ``Plt. Sekretaris MA Dorong Penguatan Infrastruktur Peradilan pada Tahun Anggaran 2027.'' \textit{Mahkamah Agung RI}, \url{https://www.mahkamahagung.go.id/id/berita/7541/plt-sekretaris-ma-dorong-penguatan-infrastruktur-peradilan-pada-tahun-anggaran-2027}.
\end{hangparas}

\vspace{0.15em}
{\raggedright\noindent \textbf{Peraturan Perundang-undangan dan Putusan Peradilan}\par}
\begin{hangparas}{0.4cm}{1}
    Undang-Undang Dasar Negara Republik Indonesia Tahun 1945.

    Constituci\'{o}n Pol\'{\i}tica de la Rep\'{u}blica de Costa Rica (1949).

    The 1987 Constitution of the Republic of the Philippines.

    United States Code, Title 31, Section 1105(b) (31 U.S.C. \S\ 1105(b)).

    Undang-Undang Nomor 14 Tahun 1985 tentang Mahkamah Agung sebagaimana telah diubah terakhir dengan Undang-Undang Nomor 3 Tahun 2009.

    Undang-Undang Nomor 35 Tahun 1999 tentang Perubahan atas Undang-Undang Nomor 14 Tahun 1970 tentang Ketentuan-Ketentuan Pokok Kekuasaan Kehakiman.

    Undang-Undang Nomor 17 Tahun 2003 tentang Keuangan Negara.

    Undang-Undang Nomor 1 Tahun 2004 tentang Perbendaharaan Negara.

    Undang-Undang Nomor 4 Tahun 2004 tentang Kekuasaan Kehakiman.

    Undang-Undang Nomor 48 Tahun 2009 tentang Kekuasaan Kehakiman.

    Undang-Undang Nomor 49 Tahun 2009 tentang Perubahan Kedua atas Undang-Undang Nomor 2 Tahun 1986 tentang Peradilan Umum.

    Undang-Undang Nomor 50 Tahun 2009 tentang Perubahan Kedua atas Undang-Undang Nomor 7 Tahun 1989 tentang Peradilan Agama.

    Undang-Undang Nomor 51 Tahun 2009 tentang Perubahan Kedua atas Undang-Undang Nomor 5 Tahun 1986 tentang Peradilan Tata Usaha Negara.

    Undang-Undang Nomor 5 Tahun 2014 tentang Aparatur Sipil Negara.

    Undang-Undang Nomor 20 Tahun 2023 tentang Aparatur Sipil Negara.

    Peraturan Pemerintah Nomor 94 Tahun 2012 tentang Hak Keuangan dan Fasilitas Hakim yang Berada di Bawah Mahkamah Agung.

    Peraturan Pemerintah Nomor 42 Tahun 2025 tentang Perubahan Keempat atas Peraturan Pemerintah Nomor 94 Tahun 2012 tentang Hak Keuangan dan Fasilitas Hakim yang Berada di Bawah Mahkamah Agung.

    Peraturan Mahkamah Agung Republik Indonesia Nomor 7 Tahun 2015 tentang Organisasi dan Tata Kerja Kepaniteraan dan Kesekretariatan Peradilan.

    Peraturan Mahkamah Agung Republik Indonesia Nomor 5 Tahun 2025 tentang Pengadaan Hakim Pengadilan Tingkat Pertama.

    Surat Edaran Direktur Jenderal Perbendaharaan Kementerian Keuangan Republik Indonesia Nomor SE-1/PB/2026 tentang Perubahan Ketiga atas Surat Edaran Nomor SE-38/PB/2012.

    Mahkamah Konstitusi Republik Indonesia. \textit{Putusan Nomor 32/PUU-XII/2014}, perihal Pengujian Undang-Undang Nomor 5 Tahun 2014 tentang Aparatur Sipil Negara terhadap UUD 1945, diucapkan tanggal 20 April 2015.

    Mahkamah Konstitusi Republik Indonesia. \textit{Putusan Nomor 43/PUU-XIII/2015}, perihal Pengujian UU Peradilan Umum, UU Peradilan Agama, dan UU PTUN terhadap UUD 1945, diucapkan tanggal 7 Oktober 2015.

    Mahkamah Konstitusi Republik Indonesia. \textit{Putusan Nomor 66/PUU-XXI/2023}, perihal Pengujian Undang-Undang Nomor 12 Tahun 2011 tentang Pembentukan Peraturan Perundang-undangan terhadap UUD 1945.

    Mahkamah Konstitusi Republik Indonesia. \textit{Putusan Nomor 121/PUU-XXII/2024}, perihal Pengujian Undang-Undang Nomor 20 Tahun 2023 tentang Aparatur Sipil Negara terhadap UUD 1945.
\end{hangparas}

\end{document}
